\documentclass[11pt]{article}
\usepackage[margin=1in]{geometry}
\usepackage{booktabs,longtable,array,xcolor,amsmath,amssymb}
\usepackage[colorlinks=true,linkcolor=blue,urlcolor=blue,citecolor=blue]{hyperref}
\definecolor{critical}{HTML}{8B1E3F}\definecolor{newgreen}{HTML}{176B3A}\definecolor{partial}{HTML}{9A6700}
\newcommand{\prior}{\textcolor{critical}{\textbf{PRIOR ART}}}
\newcommand{\novel}{\textcolor{newgreen}{\textbf{APPARENTLY NOVEL}}}
\newcommand{\partialnovel}{\textcolor{partial}{\textbf{EXTENSION / CLARIFICATION}}}
\title{\textbf{Novelty Map: QWOA, XY Mixers and Johnson-Scheme Mixer Design}\\\large Cross-analysis of manuscript results against the verified literature}
\author{Technical literature and novelty audit}\date{23 September 2026}
\begin{document}\maketitle
\begin{abstract}
This report cross-references the principal theoretical and numerical claims of
\emph{When is a QAOA mixer a quantum walk, and which walk should it be?} against the verified bibliography maintained with the manuscript. The strongest defensible novelty is not the complete-$XY$/Johnson identity by itself, but the systematic use of the Johnson association scheme as a controlled mixer-design space and the empirical finding that feasible-graph connectivity is not monotonically related to optimization performance under the tested protocol.
\end{abstract}

\section{Executive conclusion}
Constraint preservation, $XY$ mixers, Dicke-state initialization, ring-versus-complete $XY$ comparisons, graph-informed CTQW mixer design, and topology-dependent expressivity have direct antecedents \cite{Hadfield2019,Wang2020,Cook2019,He2023,Fuchs2022,Matwiejew2024,Kordonowy2026}. The complete-$XY$ fixed-Hamming-weight dynamics also has a pre-existing Johnson/token-graph interpretation \cite{FabilaMonroy2012Token,MinSeoHeo2026}. These are foundations, not headline novelty claims.

The contribution that remains most distinctive is the systematic benchmarking of
\[
A_T=\sum_{j\in T}A_j,\qquad \emptyset\neq T\subseteq\{1,\ldots,k\},
\]
with $A_j$ the Johnson distance classes. This controlled family connects the Johnson graph to the complete feasible graph. The manuscript finds that increasing degree/density does not monotonically improve performance and that the complete feasible graph is a consistently weak endpoint under the tested protocol.

\section{Claim-by-claim novelty map}
\renewcommand{\arraystretch}{1.2}
\begin{longtable}{p{0.27\textwidth}p{0.18\textwidth}p{0.45\textwidth}}
\toprule \textbf{Claim/result} & \textbf{Status} & \textbf{Cross-analysis}\\ \midrule
\endfirsthead
\toprule \textbf{Claim/result} & \textbf{Status} & \textbf{Cross-analysis}\\ \midrule
\endhead
Constraint-preserving mixers & \prior & Established by Hadfield et al. and systematic constructions \cite{Hadfield2019,Fuchs2022,Fuchs2024}.\\
$XY$ mixers for fixed cardinality & \prior & Established; ring/complete topology studied in \cite{Wang2020,Cook2019}.\\
Dicke-state initialization & \prior & Predates this work \cite{Bartschi2019,Cook2019,He2023}; portfolio/Dicke precedent in \cite{Scursulim2026}.\\
Ring-$XY$ vs complete-$XY$ & \prior & Directly compared previously \cite{Cook2019,Wang2020,He2023}.\\
Topology affects expressivity/trainability & \prior & Prior evidence and explicit DLA analysis for $XY$ topologies \cite{Kordonowy2026}.\\
Mixer-walk dictionary under explicit hypotheses & \partialnovel & QWOA already builds mixers from graphs \cite{Marsh2019,Marsh2020}; graph-to-mixer constructions and graph-informed design exist \cite{Fuchs2022,Matwiejew2024}. The reverse dictionary is useful but not universal.\\
Complete-$XY$ at weight $k$ gives $J(n,k)$ & \prior & Follows from token-graph theory \cite{FabilaMonroy2012Token}; recent QAOA work describes $XY$ dynamics in Johnson subspaces \cite{MinSeoHeo2026}.\\
Ring-$XY$ token graph and degree $2b(x)$ & \partialnovel & Token-graph construction is prior art \cite{FabilaMonroy2012Token,Carballosa2017Regularity}. The explicit cyclic-block formula was not identified as a prior quantum-optimization result in the audited set.\\
Regularity classification & \partialnovel & Token-graph regularity is established \cite{Carballosa2017Regularity}; exact theorem priority requires careful combinatorics comparison.\\
Grover equals complete-feasible QWOA up to phase/rescaling & \partialnovel & Ingredients predate this work \cite{Bartschi2020,Slate2021,Bridi2024}; explicit identity/rescaling is a useful clarification.\\
Walk graph as design variable & \prior & Explicit in graph-informed QVA/QWOA literature \cite{Matwiejew2024,Marsh2019,Marsh2020}.\\
Johnson association scheme as systematic mixer-design family & \novel & No equivalent systematic sweep was identified in the verified bibliography. This is the clearest structural novelty.\\
Sweep of all nonempty $A_T$ at tested sizes & \novel & No matched exhaustive Johnson-scheme mixer sweep was identified.\\
Performance not monotonic in degree/density & \novel & Strongest empirical result; scope must remain restricted to tested family, instances, depths, screening and optimizer.\\
Complete feasible/Grover endpoint consistently weak & \novel & Bridi and Marquezino supply a compatible prior mechanism \cite{Bridi2024}, but not this controlled Johnson-scheme benchmark.\\
Johnson universally superior to denser mixers & \textbf{NOT SUPPORTED} & Contradicted by the manuscript's $(20,5)$ limited-depth sample. Defensible claim is non-monotonicity, not sparse-is-better.\\
Spectral degeneracy/DLA explains ordering & \textbf{HYPOTHESIS} & Motivated by \cite{Bridi2026Expressivity,Kordonowy2026}, but requires direct $\mathrm{Lie}\{iH_C,iA_T\}$ calculations.\\
Intrinsic near-term circuit advantage & \textbf{NOT YET ESTABLISHED} & Current comparison mixes approximate Trotterized Johnson with exact Grover; matched-error resources are needed \cite{He2023,Awasthi2026Trotter}.\\
\bottomrule
\end{longtable}

\section{What should explicitly be acknowledged as prior art}
The revised paper should proactively state that feasibility-preserving $XY$ mixers are established \cite{Hadfield2019,Wang2020,Fuchs2022}; ring/complete $XY$ and Dicke initialization have been compared \cite{Cook2019,He2023}; complete-$XY$/Johnson has graph-theoretic and recent algorithmic antecedents \cite{FabilaMonroy2012Token,MinSeoHeo2026}; graph-informed CTQW design already exists \cite{Matwiejew2024}; and DLA-based expressivity analysis is not new as a methodology \cite{Bridi2026Expressivity,Kordonowy2026}.

\section{What appears new}
\subsection{Johnson association scheme as controlled mixer-design space}
The strongest structural contribution is to replace isolated mixer choices by a family of Johnson distance-class unions. Each $T$ specifies which Johnson distances are traversable in one mixer step. This creates a controlled design space in which topology, degree, spectrum, distance structure and implementation locality can be compared while the feasible space is fixed.

\subsection{Non-monotonicity with connectivity}
The numerical evidence rejects a simple more-connected-is-better model. Large changes in degree need not generate comparable changes in performance; the complete graph is weak, while at the largest tested size Johnson itself is not the strongest non-complete member on the limited depth sample. Thus
\[
\text{mixer quality}\neq f(\text{degree alone}).
\]

\subsection{Complete mixing as an exceptional endpoint}
Bridi and Marquezino \cite{Bridi2024} provide a plausible prior mechanism through permutation invariance/loss of landscape structure. The present manuscript contributes a controlled Johnson-scheme experiment locating the complete endpoint relative to intermediate graphs. Prior mechanism and new empirical evidence should be clearly separated.

\section{What remains hypothesis rather than result}
\subsection{Spectral/DLA mechanism}
Report distinct eigenvalues, multiplicities and gaps for every $A_T$. For tractable sizes compute or bound
\[
\dim\operatorname{Lie}\{iH_C,iA_T\}.
\]
Then test whether these descriptors explain $P_*$, $P_{\leq5\%}$ and normalized energy better than degree alone \cite{Bridi2026Expressivity,Kordonowy2026}.

\subsection{Optimizer robustness}
The manuscript reports substantial $p=1$ optimization gaps. Scursulim et al. show strong optimizer dependence in a related Dicke-state portfolio setting \cite{Scursulim2026}. Add CMA-ES or differential evolution plus an independent multistart local method.

\subsection{Circuit resources}
Repeat the exact-Grover versus Trotterized-$XY$ comparison at matched approximation error, reporting state/observable error and CNOT/depth \cite{He2023,Awasthi2026Trotter}.

\section{Recommended contribution statement}
\begin{quote}
We use an explicit mixer-walk dictionary to place several feasibility-preserving mixers in a common graph-theoretic language and then use the Johnson association scheme as a structured design space for fixed-cardinality mixers. By systematically benchmarking unions of Johnson distance classes between the Johnson and complete feasible graphs, we show that increasing feasible-state connectivity does not monotonically improve optimization performance under the tested protocol. The complete feasible graph is a consistently weak endpoint, while the relative ordering among non-complete mixers depends on problem size and depth. These results motivate spectral and Lie-algebraic properties, rather than graph density alone, as candidate descriptors of mixer quality.
\end{quote}

\section{Recommended scientific message}
\begin{quote}
\textbf{Feasible-state connectivity is an algorithm-design resource, but more connectivity is not necessarily better. Within the Johnson association scheme, mixer performance depends on the structure of the feasible-state walk rather than on graph density alone.}
\end{quote}

\section{Priority experiments before submission}
\begin{enumerate}
\item \textbf{Spectral/DLA validation:} descriptors for every $A_T$ and DLA dimension where tractable; correlate with performance.
\item \textbf{Optimizer ablation:} global/evolutionary plus independent multistart local optimization.
\item \textbf{Matched-error circuit comparison:} compare Trotterized Johnson and complete/Grover at comparable approximation error.
\item \textbf{Unscreened sensitivity:} add a $p=1$ supplement on the full unscreened population.
\item \textbf{Multiplicity-aware statistics:} paired effect sizes, bootstrap CIs and family-wise/FDR correction for formal multi-mixer tests.
\end{enumerate}

\section{Literature crosswalk by role}
\textbf{Foundations/constrained QAOA:} \cite{Farhi2014,Hadfield2019,Wang2020,Fuchs2022,Fuchs2024}.\\
\textbf{QWOA/graph-informed design:} \cite{Marsh2019,Marsh2020,Slate2021,Matwiejew2024,Qu2024,Bennett2026Nonvariational}.\\
\textbf{XY topology/initialization/implementation:} \cite{Cook2019,He2023,Kordonowy2026,Awasthi2026Trotter,Bucher2026WarmStart}.\\
\textbf{Grover/complete mixing:} \cite{Bartschi2020,Bridi2024,Slate2021,Qu2024}.\\
\textbf{Johnson/token graphs:} \cite{FabilaMonroy2012Token,Carballosa2017Regularity,MinSeoHeo2026}.\\
\textbf{Expressivity/DLA:} \cite{Bridi2026Expressivity,Kordonowy2026}.\\
\textbf{Portfolio/Dicke context:} \cite{Hodson2019,Slate2021,He2023,Qu2024,Scursulim2026}.

\section{Scope and confidence}
APPARENTLY NOVEL means that no direct prior instance was identified in the verified bibliography assembled for this manuscript and the targeted audit behind it; it is not an absolute priority claim. The most important remaining search before submission is an exhaustive query combining Johnson association scheme, distance classes, QAOA, QWOA, quantum walk optimization and mixer design.

\bibliographystyle{apsrev4-2}
\bibliography{references}
\end{document}
